\documentclass[12pt, a4paper]{article}
\usepackage[utf8]{inputenc}
\usepackage[T1]{fontenc}
\usepackage[french]{babel}
\usepackage[margin=2.5cm]{geometry}
\usepackage{amsmath, amssymb}
\usepackage{listings}
\usepackage{xcolor}
\usepackage{tcolorbox}
\tcbuselibrary{theorems, skins, breakable}
\usepackage{hyperref}
\usepackage{booktabs}
\usepackage{fancyhdr}
\usepackage{titlesec}
\usepackage{enumitem}
\usepackage{tikz}
\usetikzlibrary{shapes, arrows.meta, positioning}

% ── Couleurs ────────────────────────────────────────────────
\definecolor{suva_blue}{RGB}{30, 100, 180}
\definecolor{code_bg}{RGB}{245, 247, 250}
\definecolor{code_kw}{RGB}{0, 80, 200}
\definecolor{code_str}{RGB}{180, 40, 40}
\definecolor{code_cmt}{RGB}{80, 130, 80}
\definecolor{warn_bg}{RGB}{255, 248, 220}
\definecolor{warn_border}{RGB}{200, 150, 0}
\definecolor{success_bg}{RGB}{230, 255, 235}
\definecolor{success_border}{RGB}{40, 160, 80}
\definecolor{dark_bg}{RGB}{25, 30, 50}

% ── Listings ─────────────────────────────────────────────────
\lstdefinestyle{bash}{
  basicstyle=\ttfamily\small\color{white},
  backgroundcolor=\color{dark_bg},
  frame=single, framerule=0pt,
  breaklines=true, showstringspaces=false,
  moredelim=[is][\color{green!70}]{\$}{\$},
}
\lstdefinestyle{python}{
  language=Python,
  basicstyle=\ttfamily\small,
  keywordstyle=\color{code_kw}\bfseries,
  stringstyle=\color{code_str},
  commentstyle=\color{code_cmt}\itshape,
  backgroundcolor=\color{code_bg},
  frame=single, rulecolor=\color{gray!40},
  numbers=left, numberstyle=\tiny\color{gray},
  breaklines=true, tabsize=4,
  showstringspaces=false,
}
\lstdefinestyle{output}{
  basicstyle=\ttfamily\small\color{black},
  backgroundcolor=\color{gray!10},
  frame=single, rulecolor=\color{gray!50},
  breaklines=true,
}

% ── Boîtes ───────────────────────────────────────────────────
\newtcolorbox{cmdbox}[1]{
  colback=dark_bg, colframe=suva_blue, colupper=white,
  fonttitle=\ttfamily\bfseries\color{white}, title={\texttt{\$} #1}
}
\newtcolorbox{infobox}[1][]{
  colback=blue!5, colframe=suva_blue,
  fonttitle=\bfseries, title={Info}, #1
}
\newtcolorbox{warningbox}[1][]{
  colback=warn_bg, colframe=warn_border,
  fonttitle=\bfseries, title={Attention}, #1
}
\newtcolorbox{successbox}[1][]{
  colback=success_bg, colframe=success_border,
  fonttitle=\bfseries, title={Succès attendu}, #1
}
\newtcolorbox{errorbox}[1][]{
  colback=red!5, colframe=red!60,
  fonttitle=\bfseries, title={Erreur connue / Solution}, #1
}

% ── En-têtes ─────────────────────────────────────────────────
% ── Commande signature ──────────────────────────────────────────────
\newcommand{\auteur}{Mohamed Lamine MBENGUE}
\newcommand{\specialites}{Sécurité · DevOps · Dev Full Stack Web \& Mobile}
\newcommand{\tel}{+221 78 178 44 36}
\newcommand{\email}{mbenguemohamedlamine2022@gmail.com}
\newcommand{\portfolio}{portfolio-alamine.web.app}

\pagestyle{fancy}
\fancyhf{}
\fancyhead[L]{\color{suva_blue}\textbf{SuVa / ETHDILITHIUM}}
\fancyhead[R]{\color{gray}Guide de Test}
\fancyfoot[L]{\footnotesize\color{gray}\auteur}
\fancyfoot[C]{\thepage}
\fancyfoot[R]{\footnotesize\color{gray}\portfolio}
\renewcommand{\headrulewidth}{0.4pt}
\renewcommand{\footrulewidth}{0.3pt}

% ────────────────────────────────────────────────────────────
\begin{document}

\begin{center}
  {\color{suva_blue}\rule{\linewidth}{2pt}}\\[0.5cm]
  {\Huge\bfseries\color{suva_blue} Guide de Test du Projet}\\[0.3cm]
  {\large\color{gray} SuVa / ETHDILITHIUM --- Mise en Route Complète}\\[0.2cm]
  {\small Mars 2026}\\[0.3cm]
  {\color{suva_blue}\rule{\linewidth}{2pt}}
\end{center}

\begin{tcolorbox}[colback=suva_blue!5, colframe=suva_blue!40,
  left=6pt, right=6pt, top=4pt, bottom=4pt]
\small
\begin{tabular}{ll}
  \textbf{Auteur} & \auteur \\
  \textbf{Spécialités} & \specialites \\
  \textbf{Tél} & \tel \quad \textbf{Email} : \email \\
  \textbf{Portfolio} & \href{https://\portfolio}{\texttt{\portfolio}} \\
\end{tabular}
\end{tcolorbox}

\tableofcontents
\newpage

%═══════════════════════════════════════════════════════════════
\section{Vue d'ensemble du projet}
%═══════════════════════════════════════════════════════════════

Le projet est organisé en deux grandes parties~:

\begin{center}
\begin{tikzpicture}[
  node distance=2cm,
  block/.style={draw=suva_blue, rounded corners, fill=blue!5,
    minimum width=4.5cm, minimum height=1cm,
    text width=4.2cm, align=center, font=\small},
  arrow/.style={->, thick, suva_blue}
]
  \node[block] (py) {\textbf{Python (off-chain)}\\\footnotesize
    \texttt{python-ref/}\\
    Génération de clés, Signature};
  \node[block, right=3cm of py] (sol) {\textbf{Solidity (on-chain)}\\\footnotesize
    \texttt{src/}\\
    Vérification on-chain};
  \draw[arrow] (py) -- node[above, font=\footnotesize] {(pk, sig)} (sol);
  \node[above=0.5cm of py, font=\scriptsize\color{gray}]
    {Machine utilisateur};
  \node[above=0.5cm of sol, font=\scriptsize\color{gray}]
    {Blockchain Ethereum};
\end{tikzpicture}
\end{center}

\bigskip
Structure du dossier principal \texttt{ETHDILITHIUM-main/}~:
\begin{lstlisting}[style=bash]
ETHDILITHIUM-main/
  makefile                  <- Commandes principales
  foundry.toml              <- Config Solidity (Foundry)
  python-ref/               <- Partie Python
    makefile
    requirements.txt
    Falcon_dilithium_py/    <- Code source Python
      Falcon_dilithium/     <- Algorithme Dilithium
      modules/              <- Arithmetique matricielle
      polynomials/          <- Anneaux polynomiaux
      Falcon_keccak_prng/   <- PRNG Keccak256
      shake/                <- Wrapper SHAKE
      utilities/            <- Utilitaires
      tests/                <- Tests unitaires
      benchmarks/           <- Benchmarks perf
  src/                      <- Contrats Solidity
  test/                     <- Tests Foundry (Solidity)
  lib/                      <- Dependances Foundry
\end{lstlisting}

%═══════════════════════════════════════════════════════════════
\section{Installation de la partie Python}
%═══════════════════════════════════════════════════════════════

\subsection{Prérequis}

\begin{itemize}
  \item Python 3.10 ou supérieur (\texttt{python3 --version})
  \item \texttt{pip} disponible
  \item \texttt{git} installé
\end{itemize}

\subsection{Étape 1 — Se placer dans le bon dossier}

\begin{cmdbox}{cd ETHDILITHIUM-main/python-ref}
\begin{lstlisting}[style=bash]
cd d:/Projects/suva/falcon_dilithium/ETHDILITHIUM-main/python-ref
\end{lstlisting}
\end{cmdbox}

\subsection{Étape 2 — Créer l'environnement virtuel}

\begin{cmdbox}{Créer le venv}
\begin{lstlisting}[style=bash]
python3 -m venv myenv
\end{lstlisting}
\end{cmdbox}

\subsection{Étape 3 — Activer l'environnement}

\begin{tcolorbox}[colback=gray!5, colframe=gray!40, title=\textbf{Windows}]
\begin{lstlisting}[style=bash]
myenv\Scripts\activate
\end{lstlisting}
\end{tcolorbox}

\begin{tcolorbox}[colback=gray!5, colframe=gray!40, title=\textbf{Linux / macOS / Git Bash}]
\begin{lstlisting}[style=bash]
source myenv/bin/activate
\end{lstlisting}
\end{tcolorbox}

Ton prompt doit afficher \texttt{(myenv)} au début.

\subsection{Étape 4 — Installer les dépendances}

\begin{cmdbox}{Installer les packages}
\begin{lstlisting}[style=bash]
pip install -r requirements.txt
pip install -e "git+https://github.com/ZKNoxHQ/NTT.git@main#egg=polyntt&subdirectory=assets/pythonref/"
\end{lstlisting}
\end{cmdbox}

\begin{infobox}
Le package \texttt{polyntt} est le moteur NTT développé par ZKNox.
Il est installé directement depuis GitHub. Tu as besoin d'une connexion internet.
\end{infobox}

\subsection{Étape 5 — Vérifier l'installation}

\begin{cmdbox}{Test rapide d'import}
\begin{lstlisting}[style=bash]
python3 -c "import polyntt; print('polyntt OK')"
python3 -c "from Crypto.Cipher import AES; print('pycryptodome OK')"
\end{lstlisting}
\end{cmdbox}

\begin{successbox}
\begin{lstlisting}[style=output]
polyntt OK
pycryptodome OK
\end{lstlisting}
\end{successbox}

%═══════════════════════════════════════════════════════════════
\section{Lancer les tests Python}
%═══════════════════════════════════════════════════════════════

\subsection{Important --- le nom du package}

\begin{warningbox}
Le dossier s'appelle \texttt{Falcon\_dilithium\_py/} mais les tests importent
depuis \texttt{dilithium\_py}. Pour résoudre ça, lance les tests
\textbf{depuis le dossier} \texttt{python-ref/} en pointant vers
\texttt{Falcon\_dilithium\_py} directement.
\end{warningbox}

\subsection{Méthode recommandée --- unittest discover}

Se placer dans \texttt{python-ref/Falcon\_dilithium\_py/} puis~:

\begin{cmdbox}{Lancer TOUS les tests}
\begin{lstlisting}[style=bash]
cd Falcon_dilithium_py
python3 -m pytest tests/ -v
\end{lstlisting}
\end{cmdbox}

Ou avec unittest~:
\begin{cmdbox}{Avec unittest}
\begin{lstlisting}[style=bash]
python3 -m unittest discover -s tests -v
\end{lstlisting}
\end{cmdbox}

\subsection{Lancer un test spécifique}

\begin{cmdbox}{Un seul fichier de test}
\begin{lstlisting}[style=bash]
python3 -m pytest tests/test_dilithium.py -v
python3 -m pytest tests/test_utils.py -v
python3 -m pytest tests/test_polynomial_generic.py -v
\end{lstlisting}
\end{cmdbox}

\subsection{Description de chaque fichier de test}

\begin{center}
\begin{tabular}{lp{9cm}}
\toprule
\textbf{Fichier} & \textbf{Ce qu'il teste} \\
\midrule
\texttt{test\_dilithium.py} &
  Le flux complet : keygen → sign → verify.
  Teste aussi que la vérification échoue avec une mauvaise clé/message.
  Inclut les KAT (Known Answer Tests) NIST. \\
\addlinespace
\texttt{test\_drbg.py} &
  Le générateur pseudo-aléatoire AES256-CTR.
  Vérifie la déterminisme : même seed → mêmes bytes. \\
\addlinespace
\texttt{test\_polynomial\_generic.py} &
  Les opérations sur les polynômes dans $R_q$.
  Addition, soustraction, multiplication via NTT. \\
\addlinespace
\texttt{test\_module\_generic.py} &
  La multiplication matrice-vecteur $A \cdot s_1$.
  Operations sur les modules (vecteurs de polynômes). \\
\addlinespace
\texttt{test\_shake.py} &
  Le wrapper SHAKE-128 et SHAKE-256. \\
\addlinespace
\texttt{test\_utils.py} &
  Les fonctions utilitaires : \texttt{decompose},
  \texttt{high\_bits}, \texttt{make\_hint}, \texttt{use\_hint}. \\
\bottomrule
\end{tabular}
\end{center}

\subsection{Sortie attendue d'un test réussi}

\begin{successbox}
\begin{lstlisting}[style=output]
test_dilithium2 ... ok
test_dilithium2_keccak_prng ... ok
test_ethdilithium2_keccak_prng ... ok
test_ethdilithium2_shake ... ok

----------------------------------------------------------------------
Ran 4 tests in 12.453s

OK
\end{lstlisting}
\end{successbox}

%═══════════════════════════════════════════════════════════════
\section{Lancer les exemples}
%═══════════════════════════════════════════════════════════════

\subsection{Exemple de base}

\begin{cmdbox}{Exemple standard Dilithium}
\begin{lstlisting}[style=bash]
cd python-ref/Falcon_dilithium_py
python3 example.py
\end{lstlisting}
\end{cmdbox}

Ce script fait~:
\begin{lstlisting}[style=python]
from Falcon_dilithium import ETHDilithium2, Dilithium2

msg = b"We are ZKNox."

# Dilithium standard
pk, sk = Dilithium2.keygen()
sig = Dilithium2.sign(sk, msg)
assert Dilithium2.verify(pk, msg, sig)
print("Dilithium2 OK")

# ETHDilithium (optimise Keccak)
pk, sk = ETHDilithium2.keygen()
sig = ETHDilithium2.sign(sk, msg)
assert ETHDilithium2.verify(pk, msg, sig)
print("ETHDilithium2 OK")
\end{lstlisting}

\subsection{Exemple ZK (BabyBear / KoalaBear)}

\begin{cmdbox}{Exemple ZK}
\begin{lstlisting}[style=bash]
python3 example_zk.py
\end{lstlisting}
\end{cmdbox}

Ce script teste les variantes ZK-friendly avec des corps alternatifs.

%═══════════════════════════════════════════════════════════════
\section{Test manuel pas-à-pas en Python interactif}
%═══════════════════════════════════════════════════════════════

Tu peux tester chaque composant séparément depuis le REPL Python.
C'est la meilleure façon de comprendre ce qui se passe.

\subsection{Test de l'arithmétique polynomiale}

\begin{lstlisting}[style=python]
# Depuis python-ref/Falcon_dilithium_py/
import sys
sys.path.insert(0, '.')

from polynomials.polynomials import PolynomialRingDilithium

# Creer l'anneau R_q = Z_q[x]/(x^256 + 1)
R = PolynomialRingDilithium()

# Creer deux polynomes aleatoires
p1 = R.random_element()  # coefficients aleatoires dans Z_q
p2 = R.random_element()

# Tester les operations
p3 = p1 + p2      # addition
p4 = p1 * p2      # multiplication (via NTT)
p5 = p1 - p2      # soustraction

print(f"Coefficients de p1[0:5] = {p1.coeffs[:5]}")
print(f"Coefficients de p3[0:5] = {p3.coeffs[:5]}")
print(f"Multiplication OK ? {p4.coeffs[0] != 0}")
\end{lstlisting}

\subsection{Test du keygen complet}

\begin{lstlisting}[style=python]
from Falcon_dilithium.Falcon_dilithium import Dilithium

# Parametres Dilithium2
dil = Dilithium(...)  # voir Falcon_default_parameters.py

# Generer une paire de cles
pk, sk = dil.keygen()
print(f"Taille cle publique : {len(pk)} bytes")
# Attendu : 1312 bytes pour Dilithium2

print(f"Taille cle secrete : {len(sk)} bytes")
# Attendu : 2528 bytes pour Dilithium2

# Signer un message
msg = b"Hello SuVa"
sig = dil.sign(sk, msg)
print(f"Taille signature : {len(sig)} bytes")
# Attendu : 2420 bytes pour Dilithium2

# Verifier
ok = dil.verify(pk, msg, sig)
print(f"Verification : {ok}")  # True
\end{lstlisting}

\subsection{Test des utilitaires}

\begin{lstlisting}[style=python]
from utilities.utils import decompose, high_bits, low_bits, make_hint, use_hint

q = 8380417
alpha = 2 * ((q - 1) // 88)  # 2 * gamma_2

# Decomposer un element
r = 1234567
r1, r0 = decompose(r, alpha)
print(f"r  = {r}")
print(f"r1 = {r1}  (bits hauts)")
print(f"r0 = {r0}  (bits bas)")
print(f"Verification : r1*alpha + r0 = {r1*alpha + r0}")

# Le hint
h = make_hint(r0, r1, alpha)
r1_recovered = use_hint(h, r, alpha)
print(f"r1 original   = {r1}")
print(f"r1 recupere   = {r1_recovered}")
print(f"Hint fonctionne : {r1 == r1_recovered}")
\end{lstlisting}

%═══════════════════════════════════════════════════════════════
\section{Benchmarks}
%═══════════════════════════════════════════════════════════════

\begin{cmdbox}{Lancer les benchmarks}
\begin{lstlisting}[style=bash]
cd python-ref/Falcon_dilithium_py
python3 -m benchmarks.benchmark_dilithium
\end{lstlisting}
\end{cmdbox}

\begin{successbox}
Résultats typiques sur une machine moderne (1000 itérations)~:
\begin{lstlisting}[style=output]
Dilithium2
  keygen  : median =  6.2 ms
  sign    : median = 27.4 ms  (moyenne 27.8 ms)
  verify  : median =  7.1 ms

Dilithium3
  keygen  : median =  9.5 ms
  sign    : median = 46.1 ms
  verify  : median = 11.3 ms
\end{lstlisting}
\end{successbox}

%═══════════════════════════════════════════════════════════════
\section{Installation et tests Solidity (Foundry)}
%═══════════════════════════════════════════════════════════════

\subsection{Installer Foundry}

\begin{cmdbox}{Installation de Foundry}
\begin{lstlisting}[style=bash]
# Installer foundryup (installer de Foundry)
curl -L https://foundry.paradigm.xyz | bash
# Puis dans un nouveau terminal :
foundryup
\end{lstlisting}
\end{cmdbox}

\begin{infobox}
Foundry installe~: \texttt{forge} (compilateur/testeur),
\texttt{cast} (interactions avec la blockchain),
\texttt{anvil} (nœud local).
Vérifie avec \texttt{forge --version}.
\end{infobox}

\subsection{Générer les vecteurs de test}

Les tests Solidity nécessitent d'abord de générer des vecteurs de test
depuis Python~:

\begin{cmdbox}{Générer les vecteurs de test}
\begin{lstlisting}[style=bash]
cd ETHDILITHIUM-main/python-ref
# Generer les vecteurs ETHDilithium
python3 -m Falcon_dilithium_py.generate_ethdilithium_test_vectors
# Generer les vecteurs Dilithium standard
python3 -m Falcon_dilithium_py.generate_dilithium_test_vectors
\end{lstlisting}
\end{cmdbox}

Cela écrit un fichier \texttt{test/ZKNOX\_dilithium.t.sol} avec les
données de test en dur pour Foundry.

\subsection{Lancer les tests Solidity}

\begin{cmdbox}{Tests Solidity rapides (profil lite)}
\begin{lstlisting}[style=bash]
cd ETHDILITHIUM-main
FOUNDRY_PROFILE=lite forge test -vv
\end{lstlisting}
\end{cmdbox}

\begin{cmdbox}{Tests Solidity complets (plus lents)}
\begin{lstlisting}[style=bash]
forge test -vv
\end{lstlisting}
\end{cmdbox}

\begin{infobox}
\texttt{-vv} = verbose level 2 (affiche les logs des tests).\\
\texttt{-vvv} = encore plus verbeux (affiche les traces d'exécution).
\end{infobox}

\subsection{Vérification gas}

Pour voir le coût en gas de chaque opération~:

\begin{cmdbox}{Rapport gas}
\begin{lstlisting}[style=bash]
forge test --gas-report
\end{lstlisting}
\end{cmdbox}

%═══════════════════════════════════════════════════════════════
\section{Erreurs fréquentes et solutions}
%═══════════════════════════════════════════════════════════════

\begin{errorbox}
\textbf{Erreur~:} \texttt{ModuleNotFoundError: No module named 'polyntt'}

\textbf{Solution~:} Le package NTT n'est pas installé.
\begin{lstlisting}[style=bash]
pip install -e "git+https://github.com/ZKNoxHQ/NTT.git@main#egg=polyntt&subdirectory=assets/pythonref/"
\end{lstlisting}
\end{errorbox}

\begin{errorbox}
\textbf{Erreur~:} \texttt{ModuleNotFoundError: No module named 'Crypto'}

\textbf{Solution~:} pycryptodome mal installé.
\begin{lstlisting}[style=bash]
pip uninstall pycryptodome pycrypto
pip install pycryptodome==3.14.1
\end{lstlisting}
\end{errorbox}

\begin{errorbox}
\textbf{Erreur~:} \texttt{ImportError: attempted relative import beyond top-level package}

\textbf{Solution~:} Lance Python depuis le bon répertoire.
Être dans \texttt{python-ref/} et utiliser \texttt{python3 -m Falcon\_dilithium\_py.tests.test\_dilithium}.
\end{errorbox}

\begin{errorbox}
\textbf{Erreur~:} \texttt{forge: command not found}

\textbf{Solution~:} Foundry n'est pas dans le PATH.
\begin{lstlisting}[style=bash]
# Ajouter ~/.foundry/bin au PATH (Linux/macOS)
export PATH="$HOME/.foundry/bin:$PATH"
# Sur Windows, redemarrer le terminal apres foundryup
\end{lstlisting}
\end{errorbox}

%═══════════════════════════════════════════════════════════════
\section{Récapitulatif des commandes clés}
%═══════════════════════════════════════════════════════════════

\begin{center}
\begin{tabular}{ll}
\toprule
\textbf{Action} & \textbf{Commande} \\
\midrule
Installer Python deps & \texttt{make install} (dans \texttt{python-ref/}) \\
Lancer tests Python & \texttt{make test} (dans \texttt{python-ref/}) \\
Lancer exemple & \texttt{make example} \\
Générer test vectors & \texttt{make generate\_test\_vectors} \\
Installer Foundry & \texttt{foundryup} \\
Tests Solidity rapides & \texttt{FOUNDRY\_PROFILE=lite forge test -vv} \\
Tests Solidity complets & \texttt{forge test -vv} \\
Rapport gas & \texttt{forge test --gas-report} \\
Compiler Solidity & \texttt{forge build} \\
\bottomrule
\end{tabular}
\end{center}

\end{document}
